\documentclass[11pt]{article}

\usepackage[margin=1in]{geometry}
\usepackage[T1]{fontenc}
\usepackage{lmodern}
\usepackage{microtype}
\usepackage{setspace}
\usepackage{amsmath,amssymb}
\usepackage{graphicx}
\usepackage{hyperref}

\title{\textbf{ICT Fair Value Gap: A Quantitative Retail Strategy Audit}}
\author{Michael Moody}
\date{\today}

\begin{document}
\maketitle

\begin{abstract}
This paper studies the ICT Fair Value Gap framework as a public retail trading strategy. The goal is to translate the strategy into explicit rules, evaluate its behavior under realistic assumptions, and determine whether any restricted version of the method appears robust after costs, slippage, and regime change.
\end{abstract}

\section{Introduction}
The ICT Fair Value Gap is widely discussed in retail trading spaces, but its implementation is often ambiguous and visually discretionary.

\section{Claim Under Examination}
This section will define the retail claim clearly.

\section{Rule Formalization}
This section will translate the strategy into explicit testable rules.

\section{Data and Backtest Design}
This section will describe data, sample period, cost assumptions, and test design.

\section{Results}
This section will report gross and net performance, parameter sensitivity, and regime dependence.

\section{Discussion}
This section will assess whether the method is misleading, fragile, or usable only under narrow constraints.

\section{Conclusion}
This section will summarize the evidence and practical implications for retail traders.

\end{document}